\documentclass[11pt,a4paper]{article}

\usepackage[utf8]{inputenc}
\usepackage[T1]{fontenc}
\usepackage{amsmath,amssymb,amsthm}
\usepackage{graphicx}
\usepackage{booktabs}
\usepackage{hyperref}
\usepackage{xcolor}
\usepackage{tcolorbox}
\usepackage[margin=25mm]{geometry}
\usepackage{xeCJK}
\setCJKmainfont{Hiragino Mincho ProN}
\setCJKsansfont{Hiragino Sans}

\newtheorem{observation}{観察}
\newtheorem{conjecture}{予想}

\title{理論補助 note: $B_2 = 1/6$ 統一仮説\\
\large 微細構造定数とダークエネルギー密度の共通 arithmetic 基盤}
\author{Wright Brothers Research Program}
\date{2026年4月}

\begin{document}
\maketitle

\begin{abstract}
Wright Brothers プロジェクトは二つの独立な物理定数を arithmetic
公式で表現してきた: (1) 微細構造定数 $\alpha^{-1} \approx 137.036$ を
Bernoulli 数・$\zeta$ 特殊値で $10^{-7}$ 精度に表す公式
~\cite{alphaArithmetic}, (2) ダークエネルギー密度パラメータ
$\Omega_\Lambda = 2\pi/9$ を Wheeler--DeWitt 時間的 DN 境界と
$|\zeta(-1)|$ から導く関係~\cite{WDW-DN-paper}. 両公式は\textbf{同じ
第二 Bernoulli 数 $B_2 = 1/6$} を central に経由する. 本 note は
この共通性の arithmetic 起源を探る: $B_2$ が単なる数値ではなく
\textbf{Spec$(\mathbb{Z})$ 上の最も基本的な構造不変量}として機能する
可能性，物理的 implications, そして両公式の統一理論の sketch を
提示する. 本書は査読論文ではなく探索 note である.
\end{abstract}

\tableofcontents

\section{動機: 二つの arithmetic 公式}

\subsection{微細構造定数の算術表現}

先行研究~\cite{alphaArithmetic}は以下の観察を報告した:
\begin{equation}\label{eq:alpha-formula}
\alpha^{-1} = \frac{2}{|B_2|} + \frac{4}{|B_4|}
+ g\left(\frac{2}{|B_2|} + \frac{4}{|B_4|}\right)
+ d(\zeta(2d{-}2) - \zeta(2d{-}1))
\end{equation}
$d = 4$ (時空次元), $B_n$ はベルヌーイ数, $g$ は補正係数 ($\sim 0.038$).

主要項を評価すると:
\begin{align}
\frac{2}{|B_2|} &= \frac{2}{1/6} = 12\\
\frac{4}{|B_4|} &= \frac{4}{1/30} = 120
\end{align}
合計 $132$. 残りの $g$ 補正と $\zeta$ 項で $137.036$ に到達.

\subsection{ダークエネルギー密度パラメータ}

先行研究~\cite{WDW-DN-paper}は Wheeler--DeWitt 時間的 DN 仮説から:
\begin{equation}\label{eq:omega-formula}
\Omega_\Lambda = \frac{2\pi}{9} \approx 0.69813
\end{equation}
を parameter-free に導出した. 鍵となる数値は 1D DN cavity の
零点エネルギー
\begin{equation}
\zeta_{\neg 2}(-1) = (1 - 2^1)\zeta(-1) = -(-1/12) = \frac{1}{12}
\end{equation}
ここで $\zeta(-1) = -B_2/2 = -1/12$ (Bernoulli の関係式).

\subsection{共通する arithmetic 要素}

両公式に共通する数学要素:

\begin{center}
\begin{tabular}{lll}
\toprule
量 & $\alpha^{-1}$ での出現 & $\Omega_\Lambda$ での出現 \\
\midrule
$B_2 = 1/6$ & $2/|B_2| = 12$ (主要項)
& $\zeta(-1) = -B_2/2 = -1/12$ \\
$B_4 = -1/30$ & $4/|B_4| = 120$
& (使用されない) \\
$\pi$ & (直接には登場せず) & $2\pi/9$ (explicit) \\
$\zeta$ 特殊値 & $\zeta(6), \zeta(7)$ & $\zeta_{\neg 2}(-1)$ \\
\bottomrule
\end{tabular}
\end{center}

\textbf{観察}: 両公式が\textbf{$B_2 = 1/6$}を central に経由する
のは偶然か必然か? これが本 note の中心問い.

\section{$B_2 = 1/6$ の数論的中心性}

\subsection{Bernoulli 数の基礎}

Bernoulli 数 $B_n$ は以下の生成関数で定義される:
\begin{equation}
\frac{x}{e^x - 1} = \sum_{n=0}^\infty B_n \frac{x^n}{n!}
\end{equation}

最初のいくつか:
$$B_0 = 1, B_1 = -1/2, B_2 = 1/6, B_3 = 0, B_4 = -1/30, B_6 = 1/42, \ldots$$

奇数次 ($n \geq 3$) は全て 0. 偶数次が non-trivial. 絶対値は
$n$ とともに急速に増大:
$$|B_2| = 1/6, |B_4| = 1/30, |B_6| = 1/42, \ldots, |B_{12}| = 691/2730, \ldots$$

実は大きな $n$ では $|B_{2n}| \sim \frac{2(2n)!}{(2\pi)^{2n}}$ で指数的増大.

\subsection{$B_2$ が最小の non-trivial 値}

偶数次 Bernoulli 数で最小の non-zero は $B_2 = 1/6$. これは
arithmetic 的に以下の意味で特別:

\begin{enumerate}
\item \textbf{最初の non-trivial な zeta 特殊値を与える}:
$\zeta(-1) = -B_2/2$ は最初の non-trivial 整数値
\item \textbf{最初の non-trivial な Casimir エネルギーを決める}:
1D parallel plate Casimir $E \propto \zeta(-1)$
\item \textbf{$\zeta(2) = \pi^2/6 = \pi^2 B_2 \cdot 6$... 正確には
$\zeta(2) = \pi^2/6 = \pi^2 \cdot (B_2)$ modulo factors}
\item \textbf{bosonic string の critical dimension}: $26 = 24 + 2$ で
$24 = (1/B_2 \cdot 4)^{}$... 正確には bosonic string の冗長度が
$\propto 1/|B_2|$.
\end{enumerate}

\subsection{6 という数の特別さ}

$1/|B_2| = 6$. 6 は多くの基本的 arithmetic 場所で現れる:

\begin{itemize}
\item $\zeta(2) = \pi^2/6$ (Basel 問題, Euler 1734)
\item 三角形の角度和 = $\pi = 6 \times (\pi/6)$
\item SU(3) 基本表現の次元 $= 3 = 6/2$
\item 正六角形 (最効率的な平面敷き詰めの 1 つ)
\item 原子番号 $6$ = 炭素 (生命の基盤)
\item $S_3$ (最小の非アーベル群) の位数 $= 6$
\end{itemize}

これらの多くは偶然だが，$\zeta(2)$ との結びつきは数学的に深い:
\begin{equation}
\zeta(2) = \sum_{n=1}^\infty \frac{1}{n^2} = \frac{\pi^2}{6}
\end{equation}
このオイラーによる evaluation は，$B_2 = 1/6$ が自然数の平方和を
$\pi$ と結びつける役割を果たすことを示す.

\subsection{Euler 積との関係}

Riemann zeta の Euler 積:
\begin{equation}
\zeta(s) = \prod_{p} \frac{1}{1 - p^{-s}}
\end{equation}

$B_2$ は\textbf{全ての素数}の寄与を含む形で $\zeta(-1)$ に現れる.
一方 Wright Brothers の $\zeta_{\neg 2}$ は $p=2$ を\textbf{除いた}
Euler 積:
\begin{equation}
\zeta_{\neg 2}(s) = (1 - 2^{-s})\zeta(s) = \prod_{p \geq 3}\frac{1}{1 - p^{-s}}
\end{equation}

両者の比:
\begin{equation}
\frac{\zeta_{\neg 2}(s)}{\zeta(s)} = 1 - 2^{-s}
\end{equation}

$s = -1$ で $1 - 2 = -1$ (sign flip), $s = -3$ で $1 - 8 = -7$.

重要な観察: \textbf{$\zeta_{\neg 2}(-1) = +B_2/2 = +1/12$} で, $B_2$ の
\textbf{正の符号}が自然に出る. これは $\alpha^{-1}$ 公式の $2/|B_2| = 12$
とも整合する (どちらも正).

\section{物理的帰結}

\subsection{$B_2 = 1/6$ は「ポテンシャル定数」}

物理的直観: $B_2$ は「最も基本的な quantum 揺らぎ量」を決める
数字である. 具体的:
\begin{itemize}
\item 1D 調和振動子の零点エネルギー: $\hbar\omega/2$.
「$1/2$」は universal だが，複数モードの零点和を取ると
$\sum n^{-s}|_{s=-1} = -1/12$ が現れる.
$-1/12 = -B_2/2$
\item 2D 平行平板 Casimir: $-\pi^2/720$ に比例.
$720 = 6! = 6 \times 5!$. $B_2 = 1/6$ が分母に.
\item 3D Casimir: $-\pi^2/720 \cdot (\text{幾何因子})$. やはり $B_2$.
\end{itemize}

すなわち「真空揺らぎの定量化」は $B_2$ に commit する.

\subsection{$\alpha^{-1}$ での $B_2$: なぜか}

$\alpha^{-1} = 2/|B_2| + \ldots = 12 + \ldots$ の導出は何を意味するか?
先行研究~\cite{alphaArithmetic}によれば:
\begin{itemize}
\item $12$ 項 $= 2/|B_2|$ は QED の leading 寄与に対応する
\item $120$ 項 $= 4/|B_4|$ は次の order
\item 両者合わせて 132, これに QED running を加えて 137.036
\end{itemize}

QED の摂動展開で現れる Bernoulli 数は，\textbf{Feynman 図の
loop 和が $\zeta$ 値を通じて表される}ことから来る. 具体的には:
vacuum polarization 振幅は $\zeta(n)$ の係数を含み，解析接続で
$\zeta(-n) \sim B_{n+1}$ が登場する.

結果: \textbf{電磁結合定数の値自身が，Bernoulli 数の arithmetic
階層を反映している}という仮説が生まれる.

\subsection{$\Omega_\Lambda$ での $B_2$: なぜか}

$\Omega_\Lambda = 2\pi/9$ で $B_2$ は $\zeta(-1) = -B_2/2 = -1/12$ として
現れる. WDW 時間的 DN 仮説~\cite{WDW-DN-paper}では:

\begin{enumerate}
\item 1D 真空 (宇宙の時間方向) の零点エネルギー $\propto \zeta(-1)$
\item Dirichlet-Neumann 境界で $\zeta \to \zeta_{\neg 2}$ truncation
\item $|\zeta_{\neg 2}(-1)| = B_2/2 = 1/12$ が係数
\item CKN holographic bound で $\rho_\Lambda = (B_2/2) \rho_P (\ell_P/\ell_H)^2$
\item 無次元化して $\Omega_\Lambda = 2\pi \times (B_2/2) \times (3 H_0^2/(8\pi G))^{-1} \cdot \ldots = 2\pi/9$
\end{enumerate}

すなわち\textbf{$B_2$ が宇宙の真空エネルギー密度を $\pi$ と結びつける}:
\begin{equation}
\Omega_\Lambda = \frac{2\pi}{9} = \frac{8\pi B_2}{3} \cdot \text{(dimensional factor)}
\end{equation}
整理すると:
\begin{equation}
\boxed{\Omega_\Lambda = \frac{2\pi}{9} = \frac{8\pi}{3} \cdot |\zeta_{\neg 2}(-1)| = \frac{8\pi}{3} \cdot \frac{|B_2|}{2} = \frac{4\pi |B_2|}{3} = \frac{4\pi}{18} = \frac{2\pi}{9}}
\end{equation}
(trivially verifiable)

\section{統一公式の sketch}

\subsection{共通の arithmetic 構造}

両公式を再整理:
\begin{align}
\alpha^{-1} &= \frac{2}{|B_2|} + \frac{4}{|B_4|} + (\text{small corrections})\\
&= 12 + 120 + \ldots = 132 + \ldots\\
\Omega_\Lambda &= \frac{8\pi}{3}|B_2|\cdot\frac{1}{2}\\
&= \frac{4\pi|B_2|}{3} = \frac{2\pi}{9}
\end{align}

両者を同じ「基数」で書くと，$|B_2|$ が\textbf{reciprocal 形式}で
$\alpha^{-1}$ に，\textbf{direct 形式}で $\Omega_\Lambda$ に入る:
\begin{itemize}
\item $\alpha^{-1}$: $|B_2|$ in denominator → large number
\item $\Omega_\Lambda$: $|B_2|$ in numerator → small number (but multiplied by $4\pi/3 \sim 4$)
\end{itemize}

これが何を意味するか? 可能性:

\subsection{仮説: $B_2$ は「中間スケール」の determinator}

物理量には\textbf{UV (近距離/高エネルギー)}と\textbf{IR (遠距離/低エネルギー)}
の 2 種類がある. $\alpha^{-1}$ は electromagnetic coupling で UV 側
(Planck スケールまで running), $\Omega_\Lambda$ は IR 側 (Hubble
スケール).

仮説: $B_2$ は両者を\textbf{holographic に結びつける}中間スケールの
定数. 具体的:
\begin{equation}
\alpha^{-1} \cdot \Omega_\Lambda \sim \frac{1}{|B_2|} \cdot |B_2| \times (\text{O(1) factor})
\end{equation}

数値検証:
\begin{align}
\alpha^{-1} \cdot \Omega_\Lambda &\approx 137.036 \times 0.6981 \approx 95.67\\
1 \cdot \frac{4\pi}{3} \cdot \left(\frac{2/|B_2| + 4/|B_4|}{|B_2|^{-1}}\right) &= ?
\end{align}

試行: $\alpha^{-1} \Omega_\Lambda = 95.67$. これは
\begin{align}
\frac{8\pi}{3} \cdot \frac{1}{2/|B_2|} &= \frac{8\pi |B_2|}{6} = \frac{4\pi |B_2|}{3} = 0.698\\
\alpha^{-1} \cdot 0.698 &= 95.67
\end{align}
trivially one number. Not useful.

\subsection{別の試行: $B_2$ による次元比較}

$|B_2| = 1/6$, $|B_4| = 1/30$. 比 $|B_4|/|B_2| = 1/5$. ここに 5 が出る.

$|B_2| \cdot |B_4| = 1/180$. $180 = 6 \cdot 30$, おなじみの数字.

$|B_2|^2 = 1/36$. $36 = 6^2$.

これらの combinations の中で $2\pi/9$ になるものはあるか?
\begin{equation}
\frac{2\pi}{9} = \frac{2\pi}{9}
\end{equation}
$9 = 3^2$. $B_2 = 1/(2 \cdot 3)$. だから $9 = 3^2 = (6|B_2|)^2/... $
煩雑だが:
\begin{equation}
\frac{1}{9} = \frac{(|B_2|)^2 \cdot 4}{1} = \frac{4}{36} = \frac{1}{9} \checkmark
\end{equation}

つまり $2\pi/9 = 8\pi \cdot |B_2|^2/2 = 4\pi |B_2|^2$ ?
\begin{equation}
4\pi |B_2|^2 = 4\pi \cdot 1/36 = \pi/9 \neq 2\pi/9
\end{equation}
半分. 正しくは $2\pi/9 = 8\pi |B_2|^2$:
\begin{equation}
8\pi |B_2|^2 = 8\pi/36 = 2\pi/9 \checkmark
\end{equation}

面白い! $\Omega_\Lambda = 8\pi |B_2|^2 = 8\pi B_2^2$.

\subsection{新しい観察}

\begin{observation}[$B_2^2$ による $\Omega_\Lambda$]
$\Omega_\Lambda = 2\pi/9 = 8\pi B_2^2$ where $B_2 = 1/6$.
等価形式:
\begin{equation}
\Omega_\Lambda = 8\pi B_2^2 = 8\pi \cdot \frac{1}{36} = \frac{2\pi}{9}
\end{equation}
\end{observation}

この形式は $|\zeta_{\neg 2}(-1)| = B_2/2$ を使う元の formulation と等価
だが，$B_2^2$ の形式は\textbf{別の解釈}を提案する可能性がある.

\subsection{$B_2^2$ の物理的意味?}

$B_2$ が 1 つの arithmetic 不変量なら，$B_2^2$ は何か?
\begin{itemize}
\item Casimir での 2 つの境界 (2 枚の板) の寄与積?
\item 真空 fluctuation の 2 乗平均 ($\langle \phi^2 \rangle$)?
\item spin structure の 2 乗 ($K_1 \otimes K_1$)?
\end{itemize}

これらは speculation だが, \textbf{$B_2^2$ が自然に現れる物理的
context} を identify できれば統一仮説の rigorous な基盤になる.

\section{拡張予測}

$B_2$ 統一仮説が正しいとすると，他の物理定数も $B_2$ (または
Bernoulli 数一般) で表現されるべき.

\subsection{ミューオン異常磁気モーメント}

先行研究~\cite{muonG2}は $a_\mu$ (ミューオン $g-2$) の観測値が
$6/|B_6| = 252$ と数値的に近いことを指摘. $B_6 = 1/42$ で $6/|B_6| = 252$.

仮説: 全ての「anomalous moment」は $1/|B_{2k}|$ に比例する.

\begin{center}
\begin{tabular}{lll}
\toprule
量 & 公式 & 数値 \\
\midrule
$\alpha^{-1}$ 主要項 & $2/|B_2|$ & 12 \\
$\alpha^{-1}$ 次項 & $4/|B_4|$ & 120 \\
$a_\mu$ 近似 & $6/|B_6|$ & 252 \\
(次は?) & $8/|B_8|$ & 240 \\
\bottomrule
\end{tabular}
\end{center}

$8/|B_8|$ = $8/(1/30) = 240$. これは何を意味するか? future prediction
の候補.

\subsection{ダークマター密度?}

現在 $\Omega_{\rm dm} \approx 0.265$. 何か $B_n$ 表現があるか?

試行:
\begin{itemize}
\item $\pi B_2 = \pi/6 = 0.524$. 違う
\item $\pi^2 B_2^2 = \pi^2/36 = 0.274$. 近い!
\item $\pi^2/36 \approx 0.2742$, 観測 $0.265$, 3\% 差
\end{itemize}

\textbf{試行的予測}: $\Omega_{\rm dm} \approx \pi^2/36 \approx 0.274$.
観測との一致は 3\%. 気になる.

Total: $\Omega_\Lambda + \Omega_{\rm dm} + \Omega_b$
\begin{align}
\Omega_\Lambda + \Omega_{\rm dm} &= \frac{2\pi}{9} + \frac{\pi^2}{36}\\
&= 0.6981 + 0.2742 = 0.9723
\end{align}
これは $\Omega_b \approx 0.028$ と flatness $\sum \Omega = 1$ を
示唆: $\Omega_b^{\rm pred} = 1 - 0.9723 = 0.0277$, 観測 $0.0493$.
大きな差.

したがって $\pi^2/36$ は $\Omega_{\rm dm}$ としては不正確. Baryon を
入れると flatness と合わない.

\textbf{honest 評価}: coincidence の可能性が高い.

\subsection{プロトン質量 / 電子質量}

$m_p/m_e \approx 1836.15$. これを Bernoulli で?

試行:
\begin{itemize}
\item $|B_2| \cdot 11016 = 1836$? 11016 = なんの数? なし
\item $|B_{12}| \cdot ... = ?$. $|B_{12}| = 691/2730$. 複雑
\end{itemize}

clean な表現は見つからない.

\section{課題と open questions}

\subsection{なぜ $B_2$ か (他の $B_{2k}$ ではない)}

$B_2$ が universally 特別な理由は?

\begin{itemize}
\item Bernoulli 数の中で non-zero の minimum
\item $\zeta(2) = \pi^2/6$ の $\pi$ への translation factor
\item 2 次元 action の leading order
\item Spin structure ($\mathbb{Z}/2$)
\end{itemize}

これらは部分的 explanation. 完全な理由は未知.

\subsection{$B_4$ は $\alpha^{-1}$ にのみ?}

$\alpha^{-1}$ は $B_2, B_4$ 両方を使うが $\Omega_\Lambda$ は $B_2$ のみ.
なぜ?

可能性:
\begin{itemize}
\item $\alpha^{-1}$ は 4 次元 QFT (loop 構造から $B_4$)
\item $\Omega_\Lambda$ は 1 次元 (minisuperspace で $B_2$ のみ)
\item 次元が異なれば異なる Bernoulli 階層
\end{itemize}

\subsection{$\pi$ の出現機序}

$\Omega_\Lambda = 2\pi/9$ の $2\pi$ はどこから?
\begin{itemize}
\item CKN bound を無次元化する factor $8\pi/3$
\item 球座標積分 $\int d\Omega = 4\pi$
\item Fourier 変換の $2\pi$
\end{itemize}

これらは物理的に natural だが，$\alpha^{-1}$ 公式で explicit な $\pi$
が出ない (小さな $\zeta$ 補正にのみ) のとは対照的.

\subsection{統一理論のスケッチ}

両公式を単一の枠組みから導く理論は?

Sketch:
\begin{enumerate}
\item 物理定数は Spec$(\mathbb{Z})$ 上の arithmetic invariants の
functional
\item $B_n, \zeta(s)$ が基本的 building blocks
\item 特定の物理量は特定の次元と境界条件で Bernoulli sum に encode
される
\item $\alpha^{-1}$: 4 次元 QED の loop 展開 → $B_2, B_4, \ldots$
\item $\Omega_\Lambda$: 1 次元 WDW + CKN → $B_2$ single
\end{enumerate}

\textbf{必要な仕事}: Loop 計算で具体的に $B_n$ 階層を再導出する. これは
既存の QED/QFT 計算から読み取れるはずだが，rigorously tied to our
framework ではまだない.

\section{honest 評価}

\subsection{強み}

\begin{enumerate}
\item $B_2$ は両方に現れる事実は観察可能で reproducible
\item $\alpha^{-1}$ 公式は既に $10^{-7}$ 精度で実験に合う (先行研究)
\item $\Omega_\Lambda = 2\pi/9$ は 2\% で観測に合う
\item 両者の\textbf{共通要素}という観察自体は non-trivial
\end{enumerate}

\subsection{弱み}

\begin{enumerate}
\item 統一公式の explicit form は未完成
\item $\Omega_{\rm dm}, m_p/m_e$ 等の拡張予測は coincidence 止まり
\item $B_2$ が central な理由の物理的根拠は gesture のレベル
\item Look-elsewhere effect の risk (多くの数を試せば何かは合う)
\end{enumerate}

\subsection{位置付け}

本 note は\textbf{探索的}であり，\textbf{validated theory}ではない.
$B_2$ 統一仮説は以下の stage:
\begin{itemize}
\item 観察 stage: ✓ (両公式に $B_2$ が現れる)
\item 仮説 stage: ✓ (統一機構の提案)
\item 理論 stage: △ (部分的 sketch)
\item 検証 stage: × (未完)
\end{itemize}

今後の発展は: (a) Loop 計算からの $B_n$ 階層の rigorous 導出,
(b) 他の物理定数の systematic な Bernoulli 表現テスト,
(c) Spec$(\mathbb{Z})$ 上の category theoretic な基盤構築.

\section{結論}

Wright Brothers プロジェクトの二つの arithmetic 公式 ---
$\alpha^{-1}$ (微細構造定数) と $\Omega_\Lambda$ (ダークエネルギー
密度) --- は両方\textbf{第二 Bernoulli 数 $B_2 = 1/6$} を central に
使う. この共通性は以下を示唆する:

\begin{enumerate}
\item Physical constants は個別の coincidence ではなく，
Spec$(\mathbb{Z})$ 上の\textbf{共通 arithmetic 基盤}から出る可能性
\item $B_2$ は最も基本的な「真空定数」であり，多様な物理 context で
reappear する
\item $\Omega_\Lambda = 8\pi B_2^2$ という新しい表現は $B_2$ の power
が宇宙論量を直接与えることを示す
\end{enumerate}

ただし，本 note は探索的で，decisive な証拠には至っていない. 次の
具体的 step は loop 計算からの Bernoulli 階層の rigorous 導出と，
他の物理定数への systematic なテスト.

もしこの仮説が確認されれば，Wright Brothers プロジェクトは「宇宙 $=$
Spec$(\mathbb{Z})$」という中心仮説を，複数の独立な物理定数の
arithmetic 導出という形で支持する具体的な実績を持つことになる.

\begin{thebibliography}{99}

\bibitem{alphaArithmetic}
Wright Brothers Research Program, ``微細構造定数の算術的表現:
ベルヌーイ数・素数ギャップ・ゼータ値による公式'' (2026),
\texttt{alpha\_arithmetic\_ja.tex}.

\bibitem{WDW-DN-paper}
Wright Brothers Research Program, ``最小超空間量子宇宙論における
時間的 Dirichlet--Neumann 真空と $\Omega_\Lambda = 2\pi/9$ の
parameter-free 導出'' (2026), \texttt{wdw\_temporal\_dn\_ja.tex}.

\bibitem{muonG2}
Wright Brothers Research Program, early observation notes on
$252 = 6/|B_6|$ and muon $g-2$.

\end{thebibliography}

\end{document}
